\chapter{Desarrollo del Trabajo}
\label{chap:desarrollo}

\section{Introducción}

En este capítulo se describe el desarrollo técnico y los resultados obtenidos durante la implementación del sistema de trading cuantitativo. Se detallan los incrementos funcionales completados, los benchmarks de rendimiento observados, los casos de uso validados y los problemas técnicos enfrentados y resueltos. La implementación aquí documentada materializa los objetivos del Capítulo \ref{chap:objetivos}, sobre la base metodológica establecida en el Capítulo \ref{chap:metodologia}, y sirve de base para la discusión final del Capítulo \ref{chap:conclusiones}.

\section{Trazabilidad de Objetivos y Evidencia}

La Tabla~\ref{tab:trazabilidad_objetivos} relaciona cada objetivo específico del Capítulo~\ref{chap:objetivos} con evidencia verificable del desarrollo.

\begin{table}[H]
    \centering
    \caption{Trazabilidad de objetivos específicos con evidencia del desarrollo}
    \label{tab:trazabilidad_objetivos}
    \begin{tabular}{p{4.0cm} p{4.5cm} p{3.2cm} p{1.5cm}}
    \hline
    \textbf{Objetivo específico} & \textbf{Evidencia} & \textbf{Métrica / valor} & \textbf{Sección} \\
    \hline
    Estudiar el estado del arte del trading cuantitativo e indicadores técnicos
    & Revisión de literatura con 9 indicadores implementados (SMA, EMA, RSI, MACD, ATR, Bollinger, z-score, ROC, retorno)
    & 9 indicadores; 21 citas bibliográficas de trading
    & §\ref{sec:indicadores}, Cap.~\ref{chap:estado_del_arte} \\
    \hline
    Analizar requisitos funcionales y no funcionales de un sistema financiero de alta concurrencia
    & Arquitectura hexagonal con separación estricta de capas; bloqueos pesimistas en escrituras de saldo; timeouts granulares por tipo de operación
    & 0 condiciones de carrera detectadas en tests; 5 niveles de timeout definidos
    & §\ref{sec:datos}, §\ref{sec:concurrencia} \\
    \hline
    Diseñar una arquitectura modular que desacople orquestación transaccional del motor de inferencia
    & Puente Java-Python con protocolo IPC versionado; patrón Facade aislando el ciclo de vida de subprocesos; 4 capas hexagonales con dependencias unidireccionales
    & 0 dependencias inversas entre capas verificadas en auditoría; latencia IPC $<$1\,ms para payloads de velas
    & §\ref{sec:ipc}, §\ref{sec:patrones}, §\ref{sec:diseno_estatico} \\
    \hline
    Desarrollar un motor de ingesta y sincronización de datos capaz de almacenar masivamente candlesticks desde APIs externas
    & Pipeline ETL incremental con descarga paralela, detección de huecos y \textit{batch insert} nativo de MySQL con semántica idempotente
    & $>$50\,000 registros en inserción bulk; 60\,\% cobertura de tests del módulo
    & §\ref{sec:etl}, Incremento 2 \\
    \hline
    Implementar un pipeline de IA paramétrico con múltiples algoritmos y validación temporal
    & Motor de entrenamiento con XGBoost, LightGBM, SVM y redes neuronales; split 80/20 cronológico; warmup period dinámico; optimización bayesiana con Optuna
    & 70\,\% cobertura de tests; distribución de clases 71\,/\,18\,/\,11\,\%
    & §\ref{sec:ml}, Incrementos 6 y 7 \\
    \hline
    Evaluar y validar la plataforma mediante simulaciones históricas con métricas de rendimiento financiero
    & Motor de backtesting con comisión 0,1\,\% y \textit{slippage} 0,03\,\%; métricas: win rate, profit factor, max drawdown, ratio de Sharpe; backtesting multianual sobre $>$1\,000 velas efectivas
    & Tiempo de ejecución típico $<$30\,s; 80\,\% cobertura de tests del motor
    & Incremento 4, Tabla~\ref{tab:test_coverage} \\
    \hline
    \end{tabular}
\end{table}

\section{Manual de usuario}
\label{sec:manual}

La interfaz principal de la plataforma es una línea de comandos interactiva. Para
iniciar la aplicación basta con ejecutar el fichero JAR compilado desde el directorio
del proyecto:

\begin{center}
    \texttt{java -jar backendBotTrading-1.0.0.jar}
\end{center}

La aplicación arranca el contexto Spring\,Boot, verifica la conexión con la base de
datos MySQL y presenta el indicador de sesión. Mientras no hay usuario autenticado,
el indicador muestra \texttt{[sin sesión]}; tras el inicio de sesión pasa a mostrar
\texttt{[<nombre>]}. El comando \texttt{ayuda} lista en cualquier momento todos los
comandos disponibles con su sintaxis; \texttt{exit} o \texttt{quit} cierran la
aplicación de forma limpia, liberando carteras bloqueadas y deteniendo los procesos
Python activos.

\paragraph{Gestión de cuentas y carteras.}
El primer paso es registrarse con \texttt{signup}, que solicita de forma interactiva
nombre de usuario y contraseña; la contraseña se almacena con hash BCrypt, de modo
que nunca se persiste en texto plano. Una vez registrado, \texttt{login} autentica la
sesión y \texttt{logout} la cierra. La plataforma utiliza el concepto de
\textit{cartera de papel}: antes de operar es necesario crear al menos una cartera
con \texttt{mkpwallet <nombre>}, que solicitará el saldo inicial en USDT que actuará
como capital virtual. El comando \texttt{lw} muestra el estado de todas las carteras
del usuario, desglosando el saldo disponible, el capital comprometido en operaciones
abiertas y el capital bloqueado por instancias en ejecución.

\paragraph{Descarga de datos históricos.}
El comando \texttt{fetch} descarga velas OHLCV desde la API pública de Binance y las
persiste en la base de datos local. Requiere especificar el marco temporal
(\texttt{-tf}, por ejemplo \texttt{1h}, \texttt{4h} o \texttt{1d}) y al menos un par
de divisas (\texttt{-coins}, por ejemplo \texttt{BTCUSDT}). El argumento opcional
\texttt{-d <días>} acota la descarga a los últimos $n$ días; si se omite, se descarga
el histórico disponible según el marco temporal. El argumento
\texttt{-{}-detect-gaps} activa la detección y relleno automático de huecos en el
histórico ya almacenado sin sobreescribir los registros existentes, lo que permite
mantener la base de datos actualizada de forma incremental con ejecuciones
sucesivas. Los indicadores técnicos que necesita cada estrategia (medias móviles,
RSI, ATR, Bandas de Bollinger, entre otros) se calculan automáticamente en tiempo de
ejecución cuando se lanza un backtesting o una sesión de trading en tiempo real, por
lo que no requieren ningún paso de preprocesado adicional por parte del usuario.

\paragraph{Simulación histórica (\textit{backtesting}).}
Antes de lanzar un backtesting conviene consultar las estrategias disponibles con
\texttt{le}, que lista los scripts del directorio de estrategias junto con su
descripción. El comando \texttt{backtest} ejecuta la simulación: solicita
interactivamente el capital a asignar, el riesgo por operación expresado como
porcentaje del capital y la comisión por operación en tanto por ciento. Si se
proporciona el argumento \texttt{-model <modelo>}, el backtesting usa las predicciones
del modelo de IA para filtrar las señales de entrada, combinando la lógica de la
estrategia clásica con la capa predictiva. Los resultados se exportan en ficheros CSV
en el directorio de resultados e incluyen las métricas de rendimiento: tasa de acierto,
factor de ganancia, caída máxima y ratio de Sharpe, junto con el detalle de cada
operación abierta y cerrada durante la simulación.

\paragraph{Módulo de inteligencia artificial.}
El comando \texttt{models} muestra el catálogo de familias de algoritmos disponibles
(XGBoost, LightGBM, bosque aleatorio y red neuronal) con sus hiperparámetros
configurables y los valores por defecto. El comando \texttt{train} lanza el pipeline
de entrenamiento supervisado: descarga y prepara el conjunto de datos, calcula los
indicadores y las etiquetas de supervisión mediante la estrategia seleccionada,
ejecuta el entrenamiento y persiste el modelo en disco junto con su metadata. El
comando \texttt{optimize} añade una fase de búsqueda bayesiana de hiperparámetros
mediante Optuna, ejecutando hasta \texttt{-n-trials} iteraciones con validación
cruzada temporal de \texttt{-cv-folds} particiones; solo persiste la configuración
cuando el \textit{composite score} supera el umbral mínimo especificado con
\texttt{-min-composite}.

\paragraph{Negociación simulada en tiempo real.}
El comando \texttt{trade} abre una instancia de \textit{paper trading} que conecta
con la API de Binance para recibir velas en tiempo real. El argumento \texttt{-v}
emplea las señales de la estrategia clásica; \texttt{-r} usa las predicciones del
modelo de IA previamente entrenado. Cada instancia tiene un identificador único que
permite gestionarla de forma independiente: \texttt{lsa} lista las instancias activas
con su estado y capital comprometido; \texttt{lsd} muestra las detenidas; \texttt{stop}
suspende la ejecución preservando la posición abierta; \texttt{start} la reanuda; y
\texttt{term} liquida la posición, calcula el resultado final y devuelve el capital no
comprometido a la cartera. Las órdenes \texttt{stop -all} y \texttt{start -all}
operan sobre todas las instancias a la vez.

La Tabla~\ref{tab:comandos_a} recoge la referencia de comandos de gestión de cuentas,
carteras, datos y backtesting. La Tabla~\ref{tab:comandos_b} recoge los comandos de
negociación en tiempo real e inteligencia artificial.

\begin{table}[H]
    \centering
    \caption{Referencia de comandos: cuenta, cartera, datos y simulación histórica}
    \label{tab:comandos_a}
    \begin{tabular}{p{2.0cm} p{6.0cm} p{4.6cm}}
    \hline
    \textbf{Comando} & \textbf{Sintaxis} & \textbf{Descripción} \\
    \hline
    \multicolumn{3}{l}{\textit{Gestión de usuarios}} \\
    \hline
    \texttt{signup}   & \texttt{signup}
                      & Registro interactivo; contraseña almacenada con BCrypt \\
    \texttt{login}    & \texttt{login}
                      & Autenticación; activa la sesión del usuario \\
    \texttt{logout}   & \texttt{logout}
                      & Cierra la sesión activa \\
    \hline
    \multicolumn{3}{l}{\textit{Carteras virtuales}} \\
    \hline
    \texttt{mkpwallet} & \texttt{mkpwallet <nombre>}
                       & Crea cartera de papel; solicita saldo inicial en USDT \\
    \texttt{lw}        & \texttt{lw}
                       & Lista carteras con saldo disponible, comprometido y bloqueado \\
    \hline
    \multicolumn{3}{l}{\textit{Datos de mercado}} \\
    \hline
    \texttt{fetch} & \texttt{fetch -tf <marco> -coins <par1> [par2...]} \newline
                     \texttt{[-d <días>] [{-}{-}detect-gaps]}
                   & Descarga histórica de velas desde Binance; \texttt{-{}-detect-gaps}
                     rellena huecos de forma incremental \\
    \hline
    \multicolumn{3}{l}{\textit{Simulación histórica}} \\
    \hline
    \texttt{le}       & \texttt{le}
                      & Lista estrategias disponibles con descripción \\
    \texttt{backtest} & \texttt{backtest -strategy <nombre> -tf <marco>} \newline
                        \texttt{-coins <par1> [par2...] [-model <modelo>]}
                      & Simulación histórica; solicita capital, riesgo y comisión;
                        \texttt{-model} activa filtrado por IA \\
    \hline
    \end{tabular}
\end{table}

\begin{table}[H]
    \centering
    \caption{Referencia de comandos: negociación en tiempo real e inteligencia artificial}
    \label{tab:comandos_b}
    \begin{tabular}{p{2.0cm} p{6.0cm} p{4.6cm}}
    \hline
    \textbf{Comando} & \textbf{Sintaxis} & \textbf{Descripción} \\
    \hline
    \multicolumn{3}{l}{\textit{Paper trading en tiempo real}} \\
    \hline
    \texttt{trade}  & \texttt{trade -v|-r -strategy <nombre>} \newline
                      \texttt{-tf <marco> -coins <par1> [par2...]}
                    & Abre instancia; \texttt{-v} usa señales clásicas,
                      \texttt{-r} usa predicciones del modelo de IA \\
    \texttt{lsa}    & \texttt{lsa}                    & Lista instancias en ejecución \\
    \texttt{lsd}    & \texttt{lsd}                    & Lista instancias detenidas \\
    \texttt{start}  & \texttt{start <ID> | start -all} & Reanuda una instancia o todas \\
    \texttt{stop}   & \texttt{stop <ID> | stop -all}   & Detiene una instancia o todas \\
    \texttt{term}   & \texttt{term <ID>}
                    & Termina la instancia, liquida la posición y libera el capital \\
    \hline
    \multicolumn{3}{l}{\textit{Inteligencia artificial}} \\
    \hline
    \texttt{models}   & \texttt{models}
                      & Catálogo de algoritmos con hiperparámetros y valores por defecto \\
    \texttt{train}    & \texttt{train -model <modelo> -tf <marco>} \newline
                        \texttt{-coins <par> [-strategy <nombre>] [-d <días>]}
                      & Entrena y persiste el modelo con la configuración por defecto \\
    \texttt{optimize} & \texttt{optimize -model <modelo> -tf <marco>} \newline
                        \texttt{-coins <par> [-strategy <nombre>] [-d <días>]} \newline
                        \texttt{[-n-trials <n>] [-cv-folds <k>] [-min-composite <\%>]}
                      & Búsqueda bayesiana de hiperparámetros con Optuna;
                        persiste la configuración que supere el umbral mínimo \\
    \hline
    \end{tabular}
\end{table}

Un flujo de trabajo típico comienza con el registro mediante \texttt{signup} y el
inicio de sesión con \texttt{login}. A continuación se crea una cartera de papel con
\texttt{mkpwallet}, asignándole el saldo inicial en USDT que se usará como capital de
operaciones. Con la cartera disponible, el siguiente paso es descargar el histórico de
precios con \texttt{fetch}, indicando el marco temporal y los pares de divisas de
interés; la primera descarga obtiene todo el histórico disponible, y las siguientes
pueden usar \texttt{-{}-detect-gaps} para actualizarlo de forma incremental sin
duplicar registros.

Con los datos descargados, \texttt{le} muestra las estrategias disponibles y
\texttt{backtest} ejecuta la primera simulación histórica. En esta fase la plataforma
calcula automáticamente los indicadores técnicos que necesita la estrategia y ejecuta
el backtesting sobre todos los datos del par seleccionado, exportando el detalle de
operaciones y las métricas de rendimiento en CSV. Si se quiere incorporar el módulo
predictivo, \texttt{train} entrena un modelo con los datos ya descargados y lo persiste
en disco; \texttt{optimize} permite mejorar el resultado lanzando una búsqueda
automática de la configuración de hiperparámetros más adecuada. Una vez disponible el
modelo, se puede relanzar el backtesting con \texttt{-model} para comparar el
rendimiento de la estrategia clásica con el de la estrategia asistida por IA. Por
último, \texttt{trade -v} o \texttt{trade -r} abre una sesión de paper trading en
tiempo real que consume velas en vivo de Binance; la instancia puede pausarse, reanudarse
y terminarse en cualquier momento con \texttt{stop}, \texttt{start} y \texttt{term}.

\section{Resultados}
\label{sec:resultados}

\subsection{Cobertura de tests automáticos}
\label{sec:tests}

La estrategia de verificación del sistema combina suites de tests unitarios en Python
y Java con una cobertura total de \textbf{466 casos de prueba} distribuidos en 20
ficheros de test. La Tabla~\ref{tab:test_summary} detalla la distribución por módulo.

\begin{table}[H]
    \centering
    \caption{Resumen de suites de tests automáticos}
    \label{tab:test_summary}
    \begin{tabular}{p{5.5cm} p{3.5cm} r p{3.0cm}}
    \hline
    \textbf{Fichero de test} & \textbf{Módulo/capa} & \textbf{Tests} & \textbf{Tecnología} \\
    \hline
    \texttt{test\_strategies.py} & Estrategias Python & 48 & pytest \\
    \texttt{test\_backtest\_engine.py} & Motor de backtesting & 41 & pytest \\
    \texttt{AITrainingServiceTest} & Entrenamiento IA & 31 & JUnit 5 + Mockito \\
    \texttt{AIOptimizationServiceTest} & Optimización IA & 31 & JUnit 5 + Mockito \\
    \texttt{FetchServiceTest} & Descarga de mercado & 27 & JUnit 5 + Mockito \\
    \texttt{AccountingServiceTest} & Contabilidad/Ledger & 24 & JUnit 5 + Mockito \\
    \texttt{test\_ipc\_protocol.py} & Protocolo IPC & 23 & pytest \\
    \texttt{test\_engine\_indicators.py} & Motor de indicadores & 23 & pytest \\
    \texttt{StrategyLifecycleApplicationServiceTest} & Ciclo de vida & 23 & JUnit 5 + Mockito \\
    \texttt{SignalRetryQueueServiceTest} & Cola de reintentos & 22 & JUnit 5 + Mockito \\
    \texttt{BacktestingServiceTest} & Servicio backtesting & 20 & JUnit 5 + Mockito \\
    \texttt{CsvUtilitiesTest} & Utilidades CSV & 19 & JUnit 5 + Mockito \\
    \texttt{SignalProtocolParserTest} & Parser de señales & 16 & JUnit 5 + Mockito \\
    \texttt{StrategyQueryServiceTest} & Consulta de estrategias & 16 & JUnit 5 + Mockito \\
    \texttt{StatsCsvRepositoryTest} & Repositorio de estadísticas & 16 & JUnit 5 + Mockito \\
    \texttt{StrategyCatalogServiceTest} & Catálogo de estrategias & 15 & JUnit 5 + Mockito \\
    \texttt{MarketOrchestrationServiceTest} & Orquestación de mercado & 15 & JUnit 5 + Mockito \\
    \texttt{PaperTradingServiceTest} & Paper trading & 14 & JUnit 5 + Mockito \\
    \texttt{IndicatorsServiceTest} & Servicio indicadores & 12 & JUnit 5 + Mockito \\
    \texttt{BacktestArtifactsCleanerTest} & Limpieza de artefactos & 10 & JUnit 5 + Mockito \\
    \hline
    \textbf{Total} & & \textbf{466} & \\
    \hline
    \end{tabular}
\end{table}

Los cuatro ficheros de test Python verifican la corrección aritmética y el comportamiento
observable de los algoritmos cuantitativos con independencia del orquestador Java,
empleando datos sintéticos de semilla fija para garantizar reproducibilidad determinista.
\texttt{test\_backtest\_engine.py} cubre las funciones del motor de simulación desde el
cálculo de tamaño de posición hasta la escritura de operaciones en CSV;
\texttt{test\_engine\_indicators.py} valida la adición de columnas y su serialización;
\texttt{test\_ipc\_protocol.py} comprueba los tres modos de lectura del protocolo
MessagePack con \textit{framing} de cuatro bytes, incluyendo compatibilidad con el
formato heredado; y \texttt{test\_strategies.py} verifica de forma exhaustiva el
contrato de \texttt{BaseStrategy} y la implementación concreta
\texttt{StressTestStrategy}, desde la generación de etiquetas hasta la ausencia de
valores inválidos tras el relleno defensivo.

Los dieciséis ficheros de test Java emplean JUnit\,5 y Mockito con clases
\texttt{@Nested} que agrupan caminos felices, condiciones de error y casos límite,
sustituyendo los colaboradores externos por \textit{mocks} para aislar la unidad bajo
prueba. Las capas de mayor densidad son la de entrenamiento e inteligencia artificial
(62 tests), las estrategias (54 tests), los servicios de mercado (54 tests) y el
módulo de backtesting Java (45 tests). La integridad contable del \textit{ledger} y
las transiciones de estado del ciclo de vida de las instancias reciben cobertura
especial dada su naturaleza transaccional; el módulo de infraestructura IPC combina
16 tests de parseo de señales y 22 de la cola de reintentos para verificar el
comportamiento del puente Java-Python ante fallos y mensajes malformados.

\begin{table}[H]
    \centering
    \caption{Cobertura de tests por módulo funcional}
    \label{tab:test_coverage}
    \begin{tabular}{p{4.8cm} p{2.8cm} r r}
    \hline
    \textbf{Módulo} & \textbf{Incremento} & \textbf{Tests} & \textbf{Cobertura estimada} \\
    \hline
    Contabilidad y paper trading & Inc.\,1 y 5 & 36 & 82\,\% \\
    Pipeline de descarga y mercado & Inc.\,2 & 54 & 60\,\% \\
    Motor de indicadores Python & Inc.\,3 & 23 & 90\,\% \\
    Estrategias Python & Inc.\,3 & 48 & 90\,\% \\
    Motor de backtesting Python & Inc.\,4 & 41 & 80\,\% \\
    Backtesting Java + artefactos & Inc.\,4 & 45 & 75\,\% \\
    Protocolo IPC (Python + Java) & Inc.\,5 & 61 & 85\,\% \\
    Servicios de estrategia (Java) & Inc.\,5 & 54 & 78\,\% \\
    Entrenamiento IA & Inc.\,6 & 31 & 70\,\% \\
    Optimización IA & Inc.\,7 & 31 & 70\,\% \\
    \hline
    \textbf{Total} & & \textbf{424} & \textbf{76\,\%} \\
    \hline
    \end{tabular}
\end{table}

% TODO: añadir subsección de resultados experimentales del comando optimize
% (accuracy, F1-score, composite score, hiperparámetros óptimos, par y timeframe usados)

